\documentclass[11pt,a4paper]{article}

\usepackage[utf8]{inputenc}
\usepackage[T1]{fontenc}
\usepackage{amsmath,amssymb,amsthm}
\usepackage{mathtools}
\usepackage{geometry}
\usepackage{enumitem}
\usepackage{booktabs}
\usepackage{hyperref}

\geometry{margin=2.5cm}

% ---------------------------------------------------------------------------
% Macros
% ---------------------------------------------------------------------------
\newcommand{\RR}{\mathbb{R}}
\newcommand{\NN}{\mathcal{N}}
\newcommand{\EE}{\mathbb{E}}
\newcommand{\calF}{\mathcal{F}}
\newcommand{\calL}{\mathcal{L}}

\title{Methodology Proposal:\\
Extending ETCNN to Bermuda Option Pricing}
\author{TODO}
\date{\today}

\begin{document}
\maketitle

% ---------------------------------------------------------------------------
\section{Motivation}
% ---------------------------------------------------------------------------

% TODO: Summarise why Bermuda options are the natural next step after American
% options in the ETCNN framework (Zhang, Guo, Lu 2026).

\paragraph{TODO.}

% ---------------------------------------------------------------------------
\section{Problem Formulation}
% ---------------------------------------------------------------------------

% TODO: Write the BSM linear complementarity conditions for Bermuda options
% with exercise dates $t_1 < t_2 < \ldots < t_M = T$.

\paragraph{TODO.}

% ---------------------------------------------------------------------------
\section{Exact Exercise-Date Functions}
% ---------------------------------------------------------------------------

% TODO: Describe how g1 and g2 are adapted to the piecewise-in-time setting.

\paragraph{TODO.}

% ---------------------------------------------------------------------------
\section{Network Architecture and Loss}
% ---------------------------------------------------------------------------

% TODO: ResNet structure, input normalisation, loss terms.

\paragraph{TODO.}

% ---------------------------------------------------------------------------
\section{Proposed Experiments}
% ---------------------------------------------------------------------------

\begin{enumerate}
  \item TODO
  \item TODO
\end{enumerate}

\end{document}
